% chapters/ch27_time_horizon.tex
% Chapter 27: TimeHorizonLayer — Full Specification
% Source: NRMO_v5.tex (v4 LaTeX original) Chapter 9 (line 3933) +
%         INV-2.1-6 (line 4040) + Chapter 16 component description
%         (line 4056) + Chapter 13 v2.1 Integrated Execution Flow
%         (line 3986).
% Cross-reference: vNext simulator horizon profiles (Part VIII)
%                  multi-horizon validation (Part XII §1.3 Performance).

\chapter{TimeHorizonLayer}
\label{ch:time-horizon-layer}

% =====================================================================
\section*{v7 Realignment Note (Evaluation Horizon and Over-Defense)}
\addcontentsline{toc}{section}{v7 Realignment Note}

\nrmobox{Root-cause finding}{
A short evaluation horizon is a primary cause of over-defensive collapse.
When the rollout/scoring horizon is too short, the long-run compounding of
forward growth is not yet visible, so defensive actions win on short-term
metrics and the engine wrongly concludes that caution is optimal.
}

\paragraph{Empirical evidence (civilisation domain).}
With the two-layer structure, the same civilisation state yields opposite
``best'' actions depending solely on the evaluation horizon:
at horizon~12 the best action is defensive ($g=0.10$); at horizon~100 it is
sustainable maximum growth ($g=0.30$), while extreme growth ($g=0.55$) is
revealed as ruinous. Only a long horizon credits compounding and exposes
the true optimum.

\paragraph{Design rule.}
Each domain requires a horizon long enough to credit its compounding
dynamics. Deterministic, fast-cycle domains (e.g.\ inventory) tolerate short
horizons; long-horizon domains (civilisations, compounding portfolios)
require long horizons. Using a short horizon there reproduces the historic
``poor numbers'' by silently enforcing over-defense.

\nrmobox{Document scope and source}{
\textbf{Source}: NRMO\_v5.tex (v4 LaTeX original), Chapter 9
``Time Horizon Layer'' (line 3933), reinforced by INV-2.1-6
(line 4040), Chapter 16 component description (line 4056), and
Chapter 13 v2.1 Integrated Execution Flow Step 5 (line 3986).\\[0.3em]
\textbf{Position in v2.1 stack}: Outer-shell layer, immediately
\emph{after} Meta-Governance Layer and \emph{before} the FROZEN v2.0
NRMOSystem pipeline (Step 5 of 7 in the integrated flow).\\[0.3em]
\textbf{Architectural role}: Display-only warning layer.
\textbf{Does NOT participate in admissibility judgment.}
}

% =====================================================================
\section{Purpose and Architectural Position}
\label{sec:time-horizon-purpose}

The TimeHorizonLayer is a v2.1 outer-shell layer added to the
v2.0-FROZEN NRMOSystem pipeline. Its purpose is to surface
\emph{multi-horizon irreversibility signals} to the operator
without granting any veto authority.

The motivation is structural: a decision that appears safe at the
Immediate horizon may carry compounding irreversibility at the 5- or
20-year horizon (textbook example: a leveraged position that
optimises on quarterly P\&L while degrading 20-year recovery
optionality). NRMO Core's veto operates over a single configured
horizon $T$; TimeHorizonLayer parallelises that view across five
canonical horizons.

\nrmobox{Why this is a separate layer}{
NRMO Core's \texttt{NRMO\_CORE.veto()} retains \emph{single}
authoritative judgment under invariant INV-4 (NRMO Core veto-only
principle). Allowing the layer to influence Core would create
\emph{two} sources of admissibility, breaking the invariant.
TimeHorizonLayer therefore emits a \emph{display-only} warning;
the human operator (or upstream caller) must decide whether to
revisit the decision, but Core itself does not see the flag.
}

% =====================================================================
\section{The Five Time Horizons}
\label{sec:time-horizon-five}

\begin{table}[ht]
\centering
\small
\begin{tabularx}{\linewidth}{lcX}
\toprule
\textbf{Horizon} & \textbf{Span} & \textbf{Interpretation} \\
\midrule
Immediate & current decision moment & Single-step irreversibility
(can the action be reversed inside the same session?). \\
OneYear & 1 year & Personal-cycle irreversibility (career, health,
relationship). \\
FiveYear & 5 years & Strategic irreversibility (mid-term capital
position, institutional stance). \\
TwentyYear & 20 years & Generational irreversibility (long-term
optionality of future selves and dependents). \\
FiftyYear & 50 years & Civilisation-scale irreversibility
(structural / dynastic choices). \\
\bottomrule
\end{tabularx}
\caption[Five time horizons]{The five horizons monitored by
TimeHorizonLayer. Each is evaluated independently using
horizon-conditional projections of the current state.}
\label{tab:time-horizons}
\end{table}

\paragraph{Why five (not more, not fewer).}
The five horizons are calibrated to natural human / institutional
review cadences. \emph{Immediate} captures reflex-time judgments;
\emph{OneYear} aligns with annual review; \emph{FiveYear} with
strategic-plan windows; \emph{TwentyYear} with generational
optionality; \emph{FiftyYear} with civilisation-scale governance.
Inserting additional intermediate horizons (e.g., quarter,
decade) duplicates information without adding signal; coarser
horizons miss strategic-window irreversibilities. The five-horizon
choice is empirical: it is the minimal set that captures
\emph{distinct} failure modes observed in the case study record.

% =====================================================================
\section{Output Specification}
\label{sec:time-horizon-output}

\paragraph{Single output field.}
\texttt{warn\_flag : bool}.

\paragraph{Activation rule.}
\texttt{warn\_flag} is set to \texttt{true} when irreversibility is
detected in \emph{any} of the five time horizons.

\paragraph{Per-horizon detail (optional, display only).}
For diagnostic purposes the layer may also emit a per-horizon
detail dictionary:

\begin{lstlisting}[basicstyle=\ttfamily\small,frame=single]
{
  "warn_flag": true,
  "Immediate":  False,
  "OneYear":    False,
  "FiveYear":   True,    # 5-year irreversibility detected
  "TwentyYear": True,    # 20-year amplified
  "FiftyYear":  False
}
\end{lstlisting}

\paragraph{Output channel.}
The \texttt{warn\_flag} (and optional per-horizon detail) is
\emph{only} included in the v2.1 Integrated Flow output dictionary
under the key \texttt{time\_horizon\_warning} (Chapter~13 of v4
specification, line 3986 of NRMO\_v5.tex). It is \textbf{not}
passed to \texttt{NRMO\_CORE.veto()}.

% =====================================================================
\section{Critical Constraint: \texttt{warn\_flag} is NEVER Passed
to NRMO\_CORE}
\label{sec:time-horizon-constraint}

\nrmobox{Invariant INV-2.1-6 (verbatim)}{
\centering
\textbf{Time Horizon \texttt{warn\_flag} is NOT passed to
\texttt{NRMO\_CORE}.}\\[0.3em]
(Source: NRMO\_v5.tex, Chapter 15 v2.1 Invariants, line 4042 of
the v4 LaTeX original.)
}

\subsection{Why this matters}

\begin{enumerate}
\item \textbf{Veto-only authority preservation (INV-4)}.
NRMO Core's \texttt{NRMO\_CORE.veto()} is the sole admissibility
judge under INV-4. Letting an outer-shell heuristic modify Core's
output would create a second source of admissibility, breaking the
invariant.
\item \textbf{FROZEN v2.0 pipeline preservation (INV-2.1-1)}.
The v2.0 NRMOSystem pipeline is FROZEN. Outer-shell layers are
\emph{additive only}; they cannot mutate signals that flow into
the FROZEN pipeline.
\item \textbf{Display-only feedback channel}.
The flag is human-facing. It enables the operator to recognise
``the action is admissible by Core but I should pause and reflect
on long-horizon irreversibility,'' without forcing a decision
change.
\end{enumerate}

\subsection{What is \emph{not} prohibited}

\begin{itemize}
\item \texttt{warn\_flag} \emph{may} be displayed to the operator.
\item \texttt{warn\_flag} \emph{may} cause the operator to invoke
NRMO Core a second time with revised parameters.
\item \texttt{warn\_flag} \emph{may} influence the operator's
choice among Core-admissible options.
\item Upstream callers \emph{may} log \texttt{warn\_flag} for audit.
\end{itemize}

\subsection{What is prohibited}

\begin{itemize}
\item \texttt{warn\_flag} \textbf{must not} be wired into
\texttt{NRMO\_CORE.veto()} as a parameter.
\item \texttt{warn\_flag} \textbf{must not} be wired into APCSO,
$\varepsilon$-governance, or any FROZEN v2.0 pipeline component.
\item No code path may set \texttt{NRMO\_CORE} output to ``REJECT''
on the basis of \texttt{warn\_flag} alone.
\item No silent merging of \texttt{warn\_flag} into Core's
\texttt{verdict} is permitted.
\end{itemize}

% =====================================================================
\section{Position in the v2.1 Integrated Execution Flow}
\label{sec:time-horizon-execution-position}

The v2.1 Integrated Execution Flow is the canonical 7-step pipeline
(Chapter~13 of v4 specification, line 3986 of NRMO\_v5.tex):

\begin{lstlisting}[basicstyle=\ttfamily\small,frame=single]
1. ReflexLayer.evaluate() -> HoldImmediate? -> exit
2. ModeSelector.current_mode()
3. NonErgodicMonitor.evaluate() -> SafeExit? -> exit
4. MetaGovernanceLayer.pre_check(dsi):
     CrossDomainRuin -> SafeExit? -> force Institutional, exit
     DSImonitoring -> ForceInstitutional? -> demote mode
5. TimeHorizonLayer.evaluate() -> set warn_flag (display only)
6. NRMOSystemV20.process()  [FROZEN v2.0 pipeline]
7. Output: { mode, time_horizon_warning, verdicts }
\end{lstlisting}

\paragraph{Step 5 details.}
TimeHorizonLayer is invoked \emph{after} Meta-Governance has
either let the mode pass or demoted it, and \emph{before} the
FROZEN v2.0 NRMOSystem.process() pipeline. This positioning is
critical: invoking Step 5 \emph{before} Step 4 would mean
TimeHorizon evaluates against a possibly-pre-demotion mode;
invoking it \emph{after} Step 6 would mean the warning fires after
admissibility is already decided.

\paragraph{Output structure (Step 7).}
The output dictionary has the form:

\begin{lstlisting}[basicstyle=\ttfamily\small,frame=single]
{
  "mode":                  <Reflex|Institutional|...>,
  "time_horizon_warning":  <bool, optional per-horizon detail>,
  "verdicts":              <Core verdicts>
}
\end{lstlisting}

The three top-level fields are \emph{flat siblings}.
\texttt{time\_horizon\_warning} is \emph{not} nested inside
\texttt{verdicts}; this preserves the architectural separation
between Core's authoritative judgment and the outer-shell display
warning.

% =====================================================================
\section{Relationship to Multi-Horizon Validation in $\Omega$ Full}
\label{sec:time-horizon-omega-relation}

The v5.0 $\Omega$ Full validation (Part~XII
Section~\ref{sec:methodology-v50-perf}) reports survival across
three simulation horizons: 200 / 500 / 1000 steps. These are
\emph{simulation step counts}, not the cognitive-time spans of
TimeHorizonLayer. The relationship is one of analogy, not
identity:

\begin{table}[ht]
\centering
\small
\begin{tabularx}{\linewidth}{Xl}
\toprule
\textbf{Cognitive horizon (TimeHorizonLayer)} &
\textbf{Simulation analogue (Omega Full)} \\
\midrule
Immediate & single step. \\
OneYear & 200-step horizon (short). \\
FiveYear & 500-step horizon (medium). \\
TwentyYear & 1000-step horizon (long). \\
FiftyYear & beyond simulator's empirical reach. \\
\bottomrule
\end{tabularx}
\caption[Cognitive vs simulation horizons]{Mapping between
cognitive time horizons (TimeHorizonLayer, Part~II Ch.~27) and
simulation horizons ($\Omega$ Full, Part~IX). The mapping is
\emph{conceptual}; the two systems live in different
representational domains.}
\label{tab:time-horizon-mapping}
\end{table}

\paragraph{Why the analogy is loose.}
TimeHorizonLayer evaluates a single decision against five horizons
\emph{at the moment of decision}, using horizon-conditional state
projections. $\Omega$ Full's multi-horizon validation runs
end-to-end simulations of \emph{many} decisions over the specified
horizon and reports aggregate survival. The cognitive layer warns
about an individual choice; the simulator measures policy quality.

The Long-Horizon Drift Control mechanism in $\Omega$ Full
(Section~\ref{sec:omega-drift-control}) is the closest engineering
analogue: it explicitly penalises actions that improve short-horizon
score at the cost of long-horizon optionality. In a sense,
$\lambda_{\mathrm{drift}}$ is the simulator's automated version of
TimeHorizonLayer's human-facing warning.

% =====================================================================
\section{TimeHorizonLayer is NOT in the v2.5 Engineering Codebase}
\label{sec:time-horizon-not-in-codebase}

\nrmobox{Implementation status}{
TimeHorizonLayer is part of the \textbf{cognitive specification}
(v2.0 / v2.1, Part~II of this monograph). It is \textbf{not}
implemented in the v2.5 / v5.2 simulation codebase
(Part~X, Chapter~\ref{ch:engineering-impl}).\\[0.3em]
The simulation codebase operates on the
6-dimensional civilisation state vector $(R, E, G, O, K, X)$ and
does not include the cognitive 8-tuple Situation Parameter Model
(of which \texttt{time\_horizon} is field~\#4).
}

\paragraph{Why not implemented.}
The simulation codebase is built around a horizon parameter that
is fixed per-experiment (e.g., \texttt{horizon = 200} for smoke
runs). It does not need a multi-horizon awareness layer, because
each simulation run \emph{is} its own horizon.

The cognitive specification, by contrast, addresses real-world
human decision-making in which a single choice interacts with
multiple horizons simultaneously. TimeHorizonLayer formalises
this fundamental asymmetry between simulation and human decision
contexts.

\paragraph{When implementation would matter.}
Implementation in code becomes necessary if and when the system
moves from \emph{simulation testbed} to \emph{human-in-the-loop
deployment} (Part~XII Open Items \#7). At that point the
TimeHorizonLayer.evaluate() function must be implemented to read
from the operator's actual capital, trust, health, and time
horizons, project state forward at each of the five spans, and
emit \texttt{warn\_flag}.

% =====================================================================
\section{Connections to Other Parts}
\label{sec:time-horizon-connections}

\begin{itemize}
\item INV-2.1-6 is registered in Part~III's Unified Invariant
Registry (Chapter~\ref{ch:invariants-unified}).
\item The invariant preserves NRMO Core's veto-only authority
(INV-4, Section~\ref{sec:limit-nrmo-core}).
\item The five-horizon choice maps loosely onto the simulator's
multi-horizon validation in $\Omega$ Full
(Part~IX Section~\ref{sec:omega-validation-results}; Part~XII
Section~\ref{sec:methodology-v50-perf}).
\item The cognitive 8-tuple Situation Parameter Model (Chapter 28
of this Part) uses \texttt{time\_horizon} as parameter \#4; the
TimeHorizonLayer's five-horizon span is the discretisation of that
continuous parameter.
\item The Step-5 position in the v2.1 Integrated Flow is documented
verbatim in Section~\ref{sec:v21-integrated-flow}.
\item The Long-Horizon Drift Control mechanism in $\Omega$ Full
(Section~\ref{sec:omega-drift-control}) is the engineering
analogue of TimeHorizonLayer's warning: it operationalises the
same long-horizon-irreversibility concern within the execution
engine via $\lambda_{\mathrm{drift}}$.
\item The Open Items list (Part~XII
Section~\ref{sec:methodology-v25-honest-limits}) notes that
human-in-the-loop deployment will require this layer's actual
implementation.
\end{itemize}
